% the result of experiments
\def\resultA{-52.324}
\def\resultB{-28.966}
\def\resultC{-17.603}
\def\resultD{-35.036}
\def\resultE{-33.830}
\def\resultF{-56.067}
\def\resultG{-53.791}
\def\resultH{-46.283}
\def\resultI{-29.306}
\def\resultJ{-29.215}
\def\resultK{-29.095}
\def\resultL{-29.198}
\def\resultM{-29.265}
\def\resultN{-29.630}
\def\resultO{-29.901}
\def\resultP{-36.913}
\def\resultQ{-34.960}
\def\resultR{-32.282}

% the result of experiments (with noise)
\def\resultNoiseA{-10.219}
\def\resultNoiseB{-28.966}
\def\resultNoiseC{30.637}
\def\resultNoiseD{-20.045}
\def\resultNoiseE{-32.793}
\def\resultNoiseF{-49.925}
\def\resultNoiseG{-51.806}
\def\resultNoiseH{-45.605}
\def\resultNoiseI{-28.734}
\def\resultNoiseJ{-29.176}
\def\resultNoiseK{-29.106}
\def\resultNoiseL{-29.197}
\def\resultNoiseM{-29.264}
\def\resultNoiseN{-24.179}
\def\resultNoiseO{-29.712}
\def\resultNoiseP{-32.720}
\def\resultNoiseQ{-32.303}
\def\resultNoiseR{-31.777}

\newcounter{subsubfigure}

\newcommand{\subsubcaption}[2]{%
    \backsubcaption{#1}
    \stepcounter{subsubfigure}%
    \subcaption[\alph{subfigure}-\arabic{subsubfigure}]{#2}%
    \resetsubsubcaption
}

\newcommand{\backsubcaption}[1]{%
    \addtocounter{subfigure}{-#1}%
}

\newcommand{\resetsubsubcaption}{%
    \setcounter{subsubfigure}{1}%
}

\begin{figure*}[t]   
    \centering
    \captionsetup{font=small} 

    % 行1
    \begin{subfigure}[t]{.115\textwidth}
        \includegraphics[width=\linewidth]{fig/PWL/01withoutNoise/00original.pdf}
        \vspace*{-6mm}
        % \resetsubsubcaption
        \subcaption{Original \\ MSE [dB] \\ (Case)}
    \end{subfigure}
    \renewcommand{\thesubfigure}{\alph{subfigure}-\arabic{subsubfigure}}
    \hfill
    \begin{subfigure}[t]{.115\textwidth}
        \includegraphics[width=\linewidth]{fig/PWL/01withoutNoise/01NLPD.pdf}
        \vspace*{-6mm}
        \resetsubsubcaption
        \subcaption{NLPD~\cite{chen2015discrete} \\ \resultA \\ (UN) }
    \end{subfigure}
    \hfill
    \begin{subfigure}[t]{.115\textwidth}
        \includegraphics[width=\linewidth]{fig/PWL/01withoutNoise/03AVM.pdf}
        \vspace*{-6mm}
        \subcaption{AVM~\cite{jayawant2022practical} \\ \resultC \\ (UN) }
    \end{subfigure}
    \hfill
    \begin{subfigure}[t]{.115\textwidth}
        \includegraphics[width=\linewidth]{fig/PWL/01withoutNoise/04GSSS-UN.pdf}
        \vspace*{-6mm}
        \subcaption{GSSS~\cite{hara2022gsss} \\ \resultD \\ (UN) }
    \end{subfigure}
    \hfill
    \begin{subfigure}[t]{.115\textwidth}
        \includegraphics[width=\linewidth]{fig/PWL/01withoutNoise/05SUST-UN.pdf}
        \vspace*{-6mm}
        \subcaption{SUST~\cite{hara2023graph} \\ \resultE \\ (UN)}
    \end{subfigure}
    \hfill
    \begin{subfigure}[t]{.115\textwidth}
        \includegraphics[width=\linewidth]{fig/PWL/01withoutNoise/06Design1-UN.pdf}
        \vspace*{-6mm}
        \subcaption{\textbf{Ours (i)} \\ \textbf{\resultF} \\ (UN)}
    \end{subfigure}
    \hfill
    \begin{subfigure}[t]{.115\textwidth}
        \includegraphics[width=\linewidth]{fig/PWL/01withoutNoise/07Design2-UN.pdf}
        \vspace*{-6mm}
        \subcaption{\textbf{Ours (ii)} \\ \underline{\resultG} \\ (UN)}
    \end{subfigure}
    \hfill
    \begin{subfigure}[t]{.115\textwidth}
        \includegraphics[width=\linewidth]{fig/PWL/01withoutNoise/08Design3-UN.pdf}
        \vspace*{-6mm}
        \subcaption{\textbf{Ours (iii)} \\ \resultH \\ (UN)}
    \end{subfigure}
    \hfill
        \begin{subfigure}{.024\textwidth}
        \includegraphics[width=\linewidth]{fig/PWL/01withoutNoise/colorbar.pdf}
        % 最後の画像にはサブキャプションをつけない
    \end{subfigure}

    % 行2
    \begin{subfigure}[t]{.115\textwidth}
        \hspace*{\fill}
    \end{subfigure}
    \hfill
    \begin{subfigure}[t]{.115\textwidth}
        \includegraphics[width=\linewidth]{fig/PWL/02withNoise/01NLPD_noise.pdf}
        \vspace*{-6mm}
        \subsubcaption{7}{NLPD~\cite{chen2015discrete} \\ \resultNoiseA \\ (UN+noise) }
    \end{subfigure}
    \hfill
    \begin{subfigure}[t]{.115\textwidth}
        \includegraphics[width=\linewidth]{fig/PWL/02withNoise/03AVM_noise.pdf}
        \vspace*{-6mm}
        \subsubcaption{0}{AVM~\cite{jayawant2022practical} \\ \resultNoiseC \\ (UN+noise) }
    \end{subfigure}
    \hfill
    \begin{subfigure}[t]{.115\textwidth}
        \includegraphics[width=\linewidth]{fig/PWL/02withNoise/04GSSS-UN_noise.pdf}
        \vspace*{-6mm}
        \subsubcaption{0}{GSSS~\cite{hara2022gsss} \\ \resultNoiseD \\ (UN+noise) }
    \end{subfigure}
    \hfill
    \begin{subfigure}[t]{.115\textwidth}
        \includegraphics[width=\linewidth]{fig/PWL/02withNoise/05SUST-UN_noise.pdf}
        \vspace*{-6mm}
        \subsubcaption{0}{SUST~\cite{hara2023graph} \\ \resultNoiseE \\ (UN+noise)}
    \end{subfigure}
    \hfill
    \begin{subfigure}[t]{.115\textwidth}
        \includegraphics[width=\linewidth]{fig/PWL/02withNoise/06Design1-UN_noise.pdf}
        \vspace*{-6mm}
        \subsubcaption{0}{\textbf{Ours (i)} \\ \underline{\resultNoiseF} \\ (UN+noise)}
    \end{subfigure}
    \hfill
    \begin{subfigure}[t]{.115\textwidth}
        \includegraphics[width=\linewidth]{fig/PWL/02withNoise/07Design2-UN_noise.pdf}
        \vspace*{-6mm}
        \subsubcaption{0}{\textbf{Ours (ii)} \\ \textbf{\resultNoiseG} \\ (UN+noise)}
    \end{subfigure}
    \hfill
    \begin{subfigure}[t]{.115\textwidth}
        \includegraphics[width=\linewidth]{fig/PWL/02withNoise/08Design3-UN_noise.pdf}
        \vspace*{-6mm}
        \subsubcaption{0}{\textbf{Ours (iii)} \\ \resultNoiseH \\ (UN+noise)}
    \end{subfigure}
    \hfill
        \begin{subfigure}{.024\textwidth}
        \includegraphics[width=\linewidth]{fig/PWL/01withoutNoise/colorbar.pdf}
        % 最後の画像にはサブキャプションをつけない
    \end{subfigure}

    \vspace*{-1mm}
    \caption{An example of PWL graph signals under the smoothness prior defined on a sensor graph and its sampled and recovered signals using each method under the unconstrained case with $\numv=256, \nums=16$.
    The color of each vertex indicates the magnitude of the signal value.
    UN refers the recovered signals in the unconstrained case.
    Ours (i)-(iii) refer to the Design (i)-(iii) of the proposed method described in the section~\ref{sec:propose}, respectively.
    The best and second best results in each case are highlighted in bold and with underline, respectively.}
    % \\*Note: Smaller values indicate better recovery accuracy.
\label{fig:results_PWL_UN}
\vspace*{-4mm}
\end{figure*}